\subsection{Conclusions}

This dissertation addressed the ``SF Monoculture'' --- the congestion collapse that occurs
when standard path-planning heuristics ignore LoRaWAN protocol dynamics. The central thesis
was that treating the UAV's physical position as a direct control input for the radio
protocol would enable a DRL agent to actively resolve this bottleneck. The work demonstrates
that this thesis is correct.

\subsubsection{Summary of Contributions}

\paragraph{High-Fidelity Simulation Environment.}
A custom Gymnasium-compatible environment was developed incorporating the Two-Ray Ground
Reflection model, log-normal shadowing ($\sigma=4$\,dB), EMA-based ADR latency
($\lambda=0.1$), asynchronous 10\% duty cycles, and the LoRa Capture Effect (6\,dB
threshold). This is the first open simulation framework to model the full causal chain
from UAV position through RSSI, EMA-ADR, and Capture Effect to packet delivery within a
single Gymnasium step.

\paragraph{Domain Randomisation and Curriculum Learning.}
A three-stage curriculum combining domain randomisation~\cite{tobin2017domain} and
curriculum learning~\cite{bengio2009curriculum} produced a single DQN model that
generalises across a 100$\times$ variation in operational area and a four-fold variation
in network density without retraining.

\paragraph{Fairness-Constrained Multi-Objective Reward.}
A dense reward function with AoI urgency, buffer-overflow penalties applied to
\emph{all} action types, a discovery bonus, and an energy term prevents the agent from
hovering indefinitely over a single high-rate sensor --- a failure mode seen in earlier
formulations.

\paragraph{SF-Aware Greedy Baseline.}
A novel baseline scoring candidates by a weighted combination of SF (favouring SF7--SF9)
and distance provides a stronger protocol-aware comparison point, isolating the
contribution of long-horizon planning from protocol awareness alone.

\subsubsection{Key Findings}

All three agents achieve 100\% sensor coverage in all conditions (single exception:
DQN at $N=40$, seed 1337, 97.5\%). The DQN's throughput advantage grows with sensor
count: $+2.3\%$ cumulative reward over SF-Aware Greedy at $N=10$ ($p<0.001$, $d=2.52$),
rising to $+9.0\%$ at $N=30$ ($p=0.021$, $d=0.78$) and $+7.4\%$ at $N=40$;
the energy efficiency advantage grows from $+2.3$\,B/Wh to $+17.6$\,B/Wh ($+7.5\%$).
With twenty seeds, $N=10$, $N=20$, and $N=30$ all reach statistical significance
($p<0.05$); $N=40$ shows a medium effect ($d=0.49$, $p=0.138$) not yet significant
due to extreme inter-seed layout variance. At $N=40$, Jain's Fairness Index improves
from $0.699$ (SF-Aware) to $0.725$ (DQN), $+3.7\%$. Grid-size experiments confirm
the DQN's advantage ($+2.5\%$ at $100\times100$ and $+2.0\%$ at $300\times300$,
both $p<0.001$); the policy transitions to a fairness-first mode at
$1{,}000\times1{,}000$ where throughput is statistically indistinguishable ($p=0.96$)
but Jain's Index remains superior ($0.897$ vs.\ $0.893$).

The key emergent behaviour is strategic repositioning: the DQN devotes 2.9\% of
steps to deliberate movement ($p=0.002$), reaching positions that yield 48.7\,B per
collection step versus 42.7\,B for SF-Aware Greedy (+14\%). The EMA-ADR latency
becomes a \emph{feature}: the agent maintains proximity for multiple steps to induce
SF downgrades before repositioning, enabling orthogonal-SF concurrent collection that
purely myopic greedy policies cannot plan for.

